\chapter{Introduction}
\label{ch:introduction}

\section{Context and Motivation}
\label{sec:context}

The proliferation of connected devices in healthcare has given rise to the \gls{iomt}, an ecosystem of networked medical instruments, wearable sensors, patient monitors, and clinical infrastructure that communicate over heterogeneous network protocols. The \gls{iomt} represents one of the fastest-growing segments of the broader \gls{iot} landscape, driven by the promise of remote patient monitoring, automated diagnostics, and more efficient clinical workflows~\cite{enisa2023health}.

However, this connectivity introduces a substantially enlarged attack surface. Medical devices---many of which were originally designed for isolated clinical networks---now operate in environments where they are exposed to the same classes of cyber threats that affect general-purpose computing systems, and in some cases, to domain-specific threats that exploit protocol weaknesses unique to healthcare deployments~\cite{domenech2024ensuring}. The consequences of successful intrusions in \gls{iomt} environments extend beyond data breaches; they can directly impact patient safety when devices responsible for drug infusion, vital sign monitoring, or surgical assistance are compromised.

The European Union Agency for Cybersecurity (\gls{enisa}) reported that cyberattacks targeting the health sector increased significantly between 2021 and 2023, with ransomware, data exfiltration, and denial-of-service among the most prevalent threat categories~\cite{enisa2023health}. This trend underscores the urgent need for effective \gls{idps} mechanisms tailored to the operational constraints and communication patterns of \gls{iomt} environments.

\Gls{ml}-based \glspl{ids} have emerged as a promising approach to address this challenge. Unlike signature-based systems that can only detect previously catalogued attacks, \gls{ml}-based systems learn behavioral patterns from network traffic data, enabling the detection of both known and novel attack vectors~\cite{chio2018machine}. Nevertheless, the effectiveness of these systems is fundamentally dependent on the quality, diversity, and representativeness of the datasets used for training.

Herein lies a critical gap in the current state of research: most existing \gls{iomt} and \gls{iot} intrusion detection datasets are developed independently, using different network topologies, device configurations, feature extraction methodologies, and labeling conventions. Models trained on a single dataset tend to perform well under the conditions of that specific dataset but fail to generalize when deployed in different environments or confronted with previously unseen attack patterns~\cite{domenech2025evaluating}. This phenomenon---the lack of cross-domain generalizability---represents one of the most significant barriers to the practical deployment of \gls{ml}-based \gls{ids} in real-world healthcare settings.

\section{Problem Statement}
\label{sec:problem}

The development of robust, generalizable \gls{ids} models for \gls{iomt} environments is hindered by the following interrelated problems:

\begin{enumerate}[leftmargin=*, label=\textbf{P\arabic*.}]
    \item \textbf{Dataset isolation.} Existing benchmark datasets for \gls{iomt} and \gls{iot} intrusion detection---such as CICIoMT2024~\cite{dadkhah2024ciciomt}, CICIoT2023~\cite{neto2023ciciot2023}, and IoMT-TrafficData---are developed independently, each capturing a limited subset of attack types, network protocols, and device behaviors. No single dataset provides comprehensive coverage of the threat landscape.

    \item \textbf{Feature heterogeneity.} Different datasets employ different feature extraction tools and methodologies. For instance, one dataset may use flow-level statistics extracted by Zeek, while another uses packet-level features computed with DPKT. The resulting feature spaces are structurally incompatible, making direct concatenation or comparison infeasible without substantial preprocessing.

    \item \textbf{Label inconsistency.} Attack taxonomies vary across datasets---one may label an attack as \textit{``DDoS\_TCP\_Flood''} while another records a similar pattern as \textit{``DoS''}. This absence of a unified labeling scheme prevents meaningful cross-dataset training and evaluation.

    \item \textbf{Overfitting to single domains.} Models trained on a single dataset tend to overfit to the specific network conditions, traffic distributions, and attack characteristics present in that dataset, exhibiting poor detection rates when applied to traffic from different environments~\cite{domenech2025evaluating}.
\end{enumerate}

These problems collectively limit the practical applicability of current \gls{ml}-based \gls{ids} approaches and motivate the need for a principled methodology that integrates heterogeneous data sources into a unified, coherent dataset.

\section{Research Question}
\label{sec:rq}

This thesis is guided by the following central research question:

\begin{quote}
\textit{How can the integration of multiple heterogeneous datasets improve the performance and generalizability of a machine learning model in detecting intrusions in IoMT environments?}
\end{quote}

This question encompasses both technical challenges (how to achieve meaningful integration despite structural differences) and empirical validation (whether the resulting unified dataset yields measurable improvements in model generalization).

\section{Objectives}
\label{sec:objectives}

To address the research question, this thesis defines the following objectives:

\begin{enumerate}[leftmargin=*, label=\textbf{O\arabic*.}]
    \item \textbf{Data collection and exploratory analysis.} Acquire and thoroughly analyze multiple \gls{iomt} and \gls{iot} network traffic datasets, characterizing their structure, feature spaces, attack taxonomies, class distributions, and data quality attributes.

    \item \textbf{Standardized feature extraction.} Apply the MIOTTA-NPT preprocessing tool to extract a uniform set of statistical features from the raw packet captures of each dataset, resolving feature heterogeneity at the source by ensuring all datasets share an identical feature space.

    \item \textbf{Integration pipeline design and execution.} Design and implement a systematic pipeline for dataset integration that incorporates: label harmonization across multiple granularity levels, data quality assurance, feature scaling, and dataset provenance tracking. Execute this pipeline to produce a final, unified, \gls{ml}-ready dataset.

    \item \textbf{Validation and quality assurance.} Validate the integrity and consistency of the unified dataset through statistical analysis, ensuring that the integration process preserves the informational content of the original sources while achieving structural coherence.

    \item \textbf{Cross-domain evaluation readiness.} Establish the infrastructure for rigorous cross-domain evaluation by incorporating source attribution and stratified splitting strategies that enable assessment of model generalizability across the constituent datasets.
\end{enumerate}

Additionally, should the timeline permit, an extended objective is defined:

\begin{enumerate}[leftmargin=*, label=\textbf{O\arabic*.}, start=6]
    \item \textbf{IDS model development and evaluation (extended scope).} Train and evaluate \gls{ml}-based \gls{ids} models on the unified dataset, assessing their ability to generalize across the heterogeneous data sources and comparing their performance against models trained on individual datasets.
\end{enumerate}

\section{Scope and Limitations}
\label{sec:scope}

The primary scope of this thesis encompasses the complete data engineering pipeline---from raw packet captures to a validated, \gls{ml}-ready unified dataset. The work focuses on three principal datasets:

\begin{itemize}
    \item \textbf{CICIoMT2024}~\cite{dadkhah2024ciciomt}: A multi-protocol \gls{iomt} dataset containing 18 attack types across protocols including Wi-Fi, \gls{mqtt}, and \gls{ble}.
    \item \textbf{IoMT-TrafficData}: An \gls{iomt} network traffic dataset capturing various attack categories in medical device environments.
    \item \textbf{CICIoT2023}~\cite{neto2023ciciot2023}: A large-scale general \gls{iot} dataset providing broad attack coverage across 33 attack types.
\end{itemize}

The thesis explicitly differentiates between two integration phases: a first phase focusing on the semantic integration of the two \gls{iomt}-specific datasets, and a second phase incorporating the general \gls{iot} dataset through a pipeline-oriented approach.

The following are considered outside the primary scope but may be addressed as extensions:

\begin{itemize}
    \item Training and evaluating complete \gls{ids} models (Objective~O6).
    \item Real-time deployment or performance benchmarking of \gls{ids} systems.
    \item Generation of new synthetic network traffic or adversarial samples.
\end{itemize}

\section{Methodology Overview}
\label{sec:methodology}

The methodology adopted in this thesis follows a structured, iterative approach comprising the following stages:

\begin{enumerate}
    \item \textbf{Dataset acquisition and individual exploration.} Each dataset is acquired in its original form and subjected to comprehensive exploratory data analysis, examining feature spaces, statistical properties, class distributions, data quality, and correlation structures.

    \item \textbf{Feature standardization via MIOTTA-NPT.} Raw packet captures from each dataset are processed through the MIOTTA-NPT tool using a consistent configuration (combined mode, window size of 100 packets), yielding a uniform 73-feature statistical representation for each dataset.

    \item \textbf{Cross-dataset compatibility analysis.} The MIOTTA-processed datasets are compared to verify feature alignment and identify any remaining differences in data distributions, value ranges, or quality characteristics.

    \item \textbf{Integration pipeline design.} A multi-stage pipeline is designed, encompassing label harmonization (at binary, categorical, and fine-grained levels), missing value treatment, feature scaling, source attribution, and stratified dataset splitting.

    \item \textbf{Pipeline execution and unified dataset generation.} The designed pipeline is implemented and executed, producing the final unified dataset with all quality assurance measures applied.

    \item \textbf{Validation.} The resulting dataset is validated through statistical analysis and consistency checks to ensure it meets the requirements for \gls{ml} model training.
\end{enumerate}

\section{Thesis Structure}
\label{sec:structure}

The remainder of this thesis is organized as follows:

\begin{description}
    \item[Chapter~\ref{ch:state-of-art}] reviews the state of the art in \gls{iomt} security, \gls{ids} approaches, existing datasets, and preprocessing tools relevant to this work.

    \item[Chapter~\ref{ch:solution-design}] presents the detailed specification and design of the integration methodology, including dataset selection rationale, the MIOTTA-NPT feature extraction approach, the two-phase integration strategy, and the label harmonization scheme.

    \item[Chapter~\ref{ch:execution}] describes the development and execution of the integration pipeline, presenting the exploratory analyses, the cross-dataset compatibility verification, and the step-by-step construction of the unified dataset.

    \item[Chapter~\ref{ch:experimentation}] presents the experimental evaluation: unified dataset quality assessment, \gls{ids} model training and evaluation on the preprocessed data, and a synthesis that addresses the research question in light of those empirical results.

    \item[Chapter~\ref{ch:sustainability}] analyzes the sustainability and ethical implications of this work, covering environmental, economic, and social dimensions.

    \item[Chapter~\ref{ch:conclusions}] summarizes the contributions, discusses limitations, and outlines directions for future research.
\end{description}
